% ch6_discussion.tex
\chapter{Discussion}
\label{ch:discussion}

Ce chapitre met en perspective les choix réalisés dans le portage de
StreamKM++ vers Apache Spark. L'objectif n'est pas de reprendre les résultats
du chapitre~\ref{ch:evaluation}, mais d'en dégager les principaux apports et
les limites de l'implémentation obtenue.


% ---------------------------------------------------------------------------
\section{Apports de l'implémentation}
% ---------------------------------------------------------------------------

\paragraph{Une construction du coreset adaptée au traitement de flux.}

Le principal intérêt de StreamKM++ réside dans sa capacité à construire
progressivement un résumé des observations sans conserver l'ensemble du flux
en mémoire. Lors de la construction du coreset, les points sont consommés
progressivement par le mécanisme merge-and-reduce présenté au
chapitre~\ref{ch:algorithmes}.

La quantité de données conservée dépend ainsi principalement de la taille
$m$ du coreset et du nombre de niveaux nécessaires dans la structure de
buckets, plutôt que directement du nombre total d'observations déjà
consommées. Cette propriété permet d'envisager le traitement de flux dont la
taille finale n'est pas connue à l'avance.

Il convient toutefois de distinguer la construction du coreset du protocole
expérimental utilisé dans ce projet. Pour mesurer la SSE des centres obtenus
sur les données originales, les expériences peuvent nécessiter un parcours
supplémentaire du jeu de données. Cette seconde lecture appartient à
l'évaluation du modèle et non à l'algorithme StreamKM++ lui-même.


\paragraph{Une séparation claire entre logique algorithmique et distribution.}

L'architecture du projet sépare les composants de clustering des mécanismes
propres à Spark. Les packages \texttt{streamkm.core} et
\texttt{streamkm.coreset} contiennent respectivement les opérations liées à
k-means et aux coresets, tandis que \texttt{streamkm.spark} organise leur
exécution distribuée.

Cette séparation a permis de valider les principales propriétés
algorithmiques indépendamment de Spark avant leur intégration dans le
pipeline distribué. Les mêmes classes \texttt{BucketManager},
\texttt{CoresetTree} et \texttt{KMeans} sont ainsi réutilisées dans les
versions séquentielle et distribuée.

Les résultats de validation présentés au chapitre~\ref{ch:evaluation}
montrent notamment que les versions séquentielle et Spark produisent des
coûts très proches sur les configurations étudiées, ce qui constitue un
indicateur important de la cohérence du portage.


\paragraph{Une construction naturellement distribuable par partition.}

Le partitionnement des données proposé par Spark s'intègre bien à la structure
de StreamKM++. Chaque partition construit indépendamment un coreset local,
puis ces résumés sont fusionnés pour produire le coreset global.

Cette organisation limite les échanges de données : les observations
originales restent dans leurs partitions pendant la construction locale et
seuls les résumés doivent être combinés lors de la phase de réduction.

La qualité observée lors des premières expériences montre que cette
décomposition n'entraîne pas, sur les configurations testées, de dégradation
importante par rapport à la version séquentielle.


\paragraph{Une fusion arborescente adaptée au modèle Spark.}

L'utilisation de \texttt{treeReduce} permet d'organiser la fusion des
\texttt{BucketManager} sur plusieurs niveaux au lieu de centraliser
immédiatement toutes les opérations de combinaison sur le driver.

Ce choix répond à une limite identifiée dans l'implémentation Flink de
référence, où la fusion finale des coresets constitue une étape peu
parallélisable. L'architecture Spark permet ici d'exprimer cette réduction
sous une forme plus distribuée.

Il reste cependant nécessaire de distinguer l'intérêt architectural de ce
choix de son impact réel sur les performances. Celui-ci dépend notamment du
nombre de partitions, de la taille des coresets, du coût des opérations de
fusion et des ressources matérielles disponibles. Il doit donc être évalué
expérimentalement plutôt que supposé a priori.


\paragraph{Une évaluation distribuée de la qualité.}

Le calcul de la SSE est également distribué. Les centres, dont la taille reste
faible par rapport au jeu de données, sont diffusés aux exécuteurs avec une
variable \texttt{broadcast}. Chaque partition calcule ensuite sa contribution
au coût global avant l'agrégation finale.

Cette séparation entre construction des centres et évaluation sur les données
originales est importante. Elle évite de mesurer uniquement la qualité du
clustering sur le coreset et permet de comparer les différentes méthodes sur
une métrique commune.


\paragraph{Un portage qui exploite les mécanismes propres à Spark.}

Le passage de Flink à Spark ne se limite pas à remplacer des opérateurs par
leurs équivalents syntaxiques. Plusieurs étapes du pipeline ont dû être
réorganisées selon le modèle d'exécution de Spark.

Structured Streaming fournit par exemple explicitement les frontières des
micro-batches avec \texttt{foreachBatch}, tandis que les opérations de
\texttt{broadcast} et de \texttt{treeReduce} offrent des mécanismes directs
pour diffuser les centres et fusionner les résultats intermédiaires.

Le portage a donc conduit à conserver la logique algorithmique de StreamKM++
tout en adaptant l'orchestration au modèle d'exécution du framework cible.


% ---------------------------------------------------------------------------
\section{Limites de l'implémentation}
% ---------------------------------------------------------------------------

\paragraph{Le clustering final reste exécuté sur le driver.}

Après la construction du coreset global, \texttt{KMeans.fit} est exécuté sur
le driver. Cette étape reste donc séquentielle dans l'implémentation actuelle.

Ce choix est cohérent avec le principe de StreamKM++ : le coreset étant
sensiblement plus petit que les données originales, l'étape finale doit rester
d'un coût limité. Elle peut néanmoins devenir significative lorsque $m$, $k$
ou le nombre de redémarrages augmentent.

L'ajout de ressources Spark n'accélère pas directement cette partie du
traitement. Elle constitue donc une limite potentielle de la scalabilité
globale.


\paragraph{La version Structured Streaming repose sur un état global côté driver.}

Dans l'adaptation streaming, chaque micro-batch produit d'abord un coreset
distribué, puis ce résumé est intégré dans un \texttt{BucketManager} global
conservé sur le driver.

Cette solution simplifie la gestion de l'état entre les micro-batches, mais
elle introduit également un point central dans l'architecture.

Par ailleurs, l'insertion successive de coresets déjà pondérés dans cet état
global n'est pas strictement équivalente à l'insertion directe de tous les
points bruts dans une unique structure merge-and-reduce. La version streaming
doit donc être considérée comme une adaptation pratique du principe général,
et non comme une reproduction strictement identique du traitement
séquentiel.


\paragraph{Latence liée au traitement par micro-batches.}

Dans notre implémentation, Structured Streaming est utilisé en mode
micro-batch. Les observations sont donc traitées lorsque Spark exécute le lot
auquel elles appartiennent, et non nécessairement dès leur arrivée.

Cette organisation introduit une latence liée à la durée du micro-batch et au
temps nécessaire à son traitement. Elle distingue la solution Spark d'un
traitement événement par événement tel que celui utilisé dans l'architecture
Flink de référence.

Cette limite est surtout importante pour des applications nécessitant une
réaction à très faible latence. Elle est moins critique lorsque l'objectif
principal est de maintenir périodiquement un résumé du flux.


\paragraph{Absence d'équivalent direct à Queryable State.}

La thèse exploite les mécanismes d'état de Flink pour permettre l'interrogation
de certaines informations internes du job.

L'architecture Spark retenue dans ce projet ne fournit pas de mécanisme
strictement équivalent permettant d'interroger directement le
\texttt{BucketManager} global depuis l'extérieur du traitement.

Les modèles obtenus peuvent être transmis à un sink ou à une fonction externe,
mais cela constitue un mécanisme différent de l'exposition directe de l'état
interne proposée par Flink.


\paragraph{Sensibilité au partitionnement et à l'aléatoire.}

La construction distribuée dépend du nombre de partitions. Chaque partition
construit son propre coreset à partir d'un sous-ensemble différent des données
et utilise une graine dérivée de son indice.

Deux exécutions réalisées avec le même partitionnement et les mêmes graines
peuvent être reproduites. En revanche, modifier le nombre de partitions
modifie la décomposition des données et peut donc conduire à un coreset global
différent.

Il ne s'agit pas nécessairement d'une erreur : les algorithmes de
construction des coresets et k-means++ comportent eux-mêmes des étapes
aléatoires. Cette propriété impose cependant de contrôler explicitement le
partitionnement et les graines dans le protocole expérimental.


\paragraph{L'influence de la taille du coreset n'est pas strictement monotone.}

Les expériences réalisées sur le sous-ensemble de Covertype montrent une
amélioration du coût lorsque $m$ augmente dans certaines configurations, mais
cette tendance n'est pas strictement monotone.

Ce résultat rappelle qu'une augmentation de la taille du coreset ne garantit
pas, pour une exécution donnée, une diminution systématique de la SSE. Le
processus de construction du résumé ainsi que l'initialisation k-means++
restent aléatoires.

Une caractérisation plus robuste de l'effet de $m$ nécessiterait plusieurs
répétitions avec différentes graines, puis une analyse des moyennes et des
variations observées.


\paragraph{Absence de standardisation des variables de Covertype.}

Les expériences utilisent les 54 variables de Covertype sans transformation
de leur échelle, afin de rester cohérentes avec le protocole retenu pour ce
projet.

Certaines variables numériques peuvent prendre des valeurs de plusieurs
milliers, tandis que d'autres correspondent à des indicateurs binaires.
Dans une distance euclidienne au carré, les variables de grande amplitude
peuvent donc contribuer davantage au coût.

Une standardisation modifierait la géométrie du problème et pourrait produire
des clusters différents. Elle constitue une piste d'étude intéressante, mais
n'a pas été intégrée au protocole actuel afin de conserver les mêmes données
d'entrée pour l'ensemble des méthodes comparées.


\paragraph{Absence de mécanisme d'oubli.}

Le merge-and-reduce utilisé par StreamKM++ représente l'ensemble des
observations vues depuis le début du flux. Les données anciennes continuent
donc à contribuer au coreset même lorsque la distribution du flux évolue.

Cette propriété convient lorsque la distribution des données est relativement
stable. Elle devient en revanche une limite dans un contexte de
\textit{concept drift}, où les observations récentes devraient éventuellement
avoir davantage d'influence que les anciennes.

L'ajout d'une fenêtre glissante ou d'un mécanisme de décroissance des poids
constituerait une extension possible, mais sort du périmètre de
l'implémentation étudiée.


% ---------------------------------------------------------------------------
\section{Bilan}
% ---------------------------------------------------------------------------

Les résultats obtenus montrent que le portage de StreamKM++ vers Spark est
possible sans modifier profondément les composants algorithmiques de
l'approche. La construction de coresets locaux par partition et leur fusion
permettent de conserver une qualité de clustering proche de la version
séquentielle sur les configurations étudiées.

L'intérêt principal du portage réside dans la possibilité de distribuer les
étapes les plus coûteuses sur les données originales tout en conservant un
résumé global de taille contrôlée.

La solution reste toutefois partiellement centralisée : le clustering final
et l'état utilisé entre les micro-batches résident sur le driver. Les
performances globales dépendent également du partitionnement, du paramètre
$m$ et des coûts propres au moteur Spark.

Ces éléments montrent que l'intérêt de StreamKM++ sur Spark repose sur un
compromis entre réduction des données, qualité du clustering et coût de la
distribution. Les perspectives présentées dans le chapitre suivant découlent
directement de ces limites.